\documentclass[10pt,a4paper]{amsart}
\usepackage[margin=19mm]{geometry}
\usepackage{amsmath,amssymb,mathtools}
\usepackage[scr=rsfs,cal=euler]{mathalfa}
\usepackage{xfrac}
\usepackage[unicode,hidelinks,bookmarks=false]{hyperref}
\newcommand{\ie}{i.e.\ }
\newcommand{\cf}{cf.\ }
\newcommand{\resp}{respectively\ }
\newcommand{\Subsection}{\subsection}
\providecommand{\nobreakdash}{\nobreak\hbox{-}}
\setlength{\emergencystretch}{2em}
\multlinegap0pt
\usepackage{bm}
\usepackage{bbm}
\usepackage{dsfont}
\usepackage{xspace}
\usepackage{cancel}
\usepackage{mathtools}
\usepackage{amsthm}
\makeatletter
\@ifundefined{theorem}{\newtheorem{theorem}{Theorem}}{}
\@ifundefined{lemma}{\newtheorem{lemma}[theorem]{Lemma}}{}
\makeatother
\renewcommand{\geq}{\geqslant}
\renewcommand{\leq}{\leqslant}
\title[Non-solvable torsion-free virtually solvable groups]{Non-solvable torsion-free virtually solvable groups}
\author{Jonathan A. Hillman}
\date{Source: arXiv:2302.09513v8 (CC BY 4.0)}
\begin{document}
\maketitle

\noindent\emph{Source and licence.} Non-solvable torsion-free virtually solvable groups. Version arXiv:2302.09513v8 (CC BY 4.0), distributed under the Creative Commons Attribution 4.0 International licence, as recorded in the arXiv licence metadata for this exact version and shown on the abstract page https://arxiv.org/abs/2302.09513v8. Full rights notice: licenses/arxiv-2302.09513-CC-BY-4.0.txt.\par\smallskip

\noindent\textit{Editorial scope.} Counted unit: the Theorem whose source label is \texttt{vmetabPSL(2,7)} in the article (source0.tex lines 1585--1589, source label \texttt{vmetabPSL(2,7)}), delivered together with the article's own complete original proof of it (source0.tex lines 1591--1640, one \texttt{proof} environment of 50 lines). No mathematics is written, repaired or re-derived: statement and proof are the article's own text line for line. Required context retained uncounted: overline (lines 415--421); gl(10,Q) (lines 893--898); vsolvthm (lines 974--981); split normal (lines 1397--1403); no7 (lines 1548--1553). Every definition, display and label of those lines is retained unchanged. Omitted from this package and not redistributed: 1--414, 422--892, 899--973, 982--1396, 1404--1547, 1554--1584, 1590--1590, 1641--2075. Nothing inside a retained excerpt is omitted or altered: every retained body is delivered exactly as the article's lines are written, so no omission receipt of any kind accompanies this package. Citations: the retained excerpts carry no citation command at all, so no bibliography entry of the article is transcribed and no reference is redistributed. Reference adaptation: none was needed and none was made; every reference inside the delivered unit resolves inside it, so no unresolved reference or citation is produced. Printed slips kept literal: no printed word, symbol, hypothesis or number of the counted statement or of its proof was altered. Layout and class: the article's own class is \texttt{amsart} with packages including amssymb, amscd, amsmath, amsthm, latexsym, enumerate; the wrapper is the lane's pinned \texttt{amsart} editorial wrapper at 10pt on A4, which restates the theorem-like environments the retained text needs and adds the article's own macro definitions that the retained text needs (\texttt{\textbackslash geq}, \texttt{\textbackslash leq}), with extra wrapper packages (bm, bbm, dsfont, xspace, cancel, mathtools). The delivered numbering is the unit's own. Assets: no retained span contains a figure environment, a graphics-inclusion command, a PDF-page inclusion, a table or a TikZ picture; no image file is copied into this package, the package declares no asset dependency and no image file is redistributed with it. Licence: version arXiv:2302.09513v8 of this article is distributed under the Creative Commons Attribution 4.0 International licence, as recorded in the arXiv licence metadata for this exact version and shown on the abstract page https://arxiv.org/abs/2302.09513v8. A full rights notice is delivered at licenses/arxiv-2302.09513-CC-BY-4.0.txt. Characters: every retained excerpt is pure ASCII, so the delivered engine, XeLaTeX with its default Latin Modern text font, reports no missing character.
\begin{lemma}
\label{overline}
Let $S$ be a minimal TFNS group which is not virtually abelian.
Then $S$ has normal subgroups $V\leq{U}\leq\widetilde{S}$ such that
$S/U$ is crystallographic,  $U/V$ is finite, $V/I(V)$ has positive rank
and $U/V$ acts effectively on $V/I(V)$.
\end{lemma}

\begin{lemma}
\label{gl(10,Q)}
Let $d=4$ if $H\cong{A_5}$, 
let $d=6$ if $H\cong{PSL(2,7)}$ and let $d=7$ if $H\cong{SL(2,8)}$.
If $h(U)\leq{d}$ then either $U$ is virtually abelian or $S$ is virtually nilpotent.
\end{lemma}

\begin{theorem}
\label{vsolvthm}
Let $S$ be a finitely generated, perfect group which is TFNS but not virtually nilpotent.
If $H=S/\widetilde{S}\cong{A_5}$ then $h(S)\geq9$,
if $H\cong{PSL(2,7)}$
%or $SL(2,8)$
then $h(S)\geq13$ and if $SL(2,8)$ then $h(S)\geq14$.
\end{theorem}

\begin{lemma}
\label{split normal}
Let $G$ be a crystallographic group with translation subgroup $A$ and holonomy $H$.
Suppose that the Sylow $p$-subgroup of $H$ is a cyclic subgroup $C$ which 
acts without fixed points on $A$.
Then $G$ has a subgroup isomorphic to $N_H(C)$.
\end{lemma}

\begin{theorem}
\label{no7}
If $S/\sqrt{S}\cong{PSL(2,7)}$,
$h(S/I(\sqrt{S}))=6$ and $I(\sqrt{S})\cong\mathbb{Z}^n$, for some $n$,
then $h(S)\geq15$.
\end{theorem}

\begin{theorem}
\label{vmetabPSL(2,7)}
Let $S$ be a minimal TFNS group with $S/\widetilde{S}\cong{PSL(2,7)}$.
Then $h(S)\geq14$.
\end{theorem}

\begin{proof}
There are normal subgroups $V\leq{U}<S$ such that $S/U$ is crystallographic
and $U/V$ is a finite solvable group which acts effectively on $V/I(V)$,
by Lemma \ref{overline}.
We may assume that $S$ is not virtually nilpotent,
by the observations following Theorem \ref{no7}.
We may assume also that $h(S)=6$, for otherwise $h(S)\geq14$,
by Theorem \ref{vsolvthm}.
Then $S/U$ has  a subgroup isomorphic to $M_{7,3}$, 
by Lemma \ref{split normal}.
Let $X$ be the preimage of this subgroup in $S$.
Then $X/V$ is a finite solvable group.

Suppose that $V$ is abelian and $h=h(V)<9$.
Then $X/V$ acts effectively on $V$,
for otherwise some element of $X$ with non-trivial image in $S/U$ would centralize $V$.
But any such element must map non-trivially to the simple quotient $PSL(2,7)$,
and it would follow that $V$ must be central in $S$.
In particular, $S$ would be virtually nilpotent, 
and the action of the holonomy on $V$ would be trivial.
Hence $U=V\cong\mathbb{Z}$ or $\mathbb{Z}^2$, 
since $h(S/U)=6$ and $S/U$ has holonomy $PSL(2,7)$.
But no extension of a finite nonabelian group by $\mathbb{Z}$ or $\mathbb{Z}^2$
is torsion free.
Hence we may assume that $X/V$ embeds in $GL(8,\mathbb{Q})$.

Since $X/V$ is finite it preserves a lattice $L\leq{V}$.
Hence $X$ has a finitely generated subgroup $W$ which is an extension of $X/V$
by $L\cong\mathbb{Z}^h$.
If $W$ were torsion free then this subgroup would be a Bieberbach group 
of dimension $h\leq8$, and with holonomy of order a multiple of 21,
since $X/V$ acts effectively and maps onto $M_{7,3}$.
But there are no such Bieberbach groups. 
Hence either $I(V)\not=1$ or $h(S)\geq15$.

Thus we may assume that $h(S)=13$ and $I(V)\not=1$.
The argument for Lemma \ref{gl(10,Q)} may be used to show 
that $\sqrt{S}$ is abelian of rank 6.
Hence $h(V/I(V))=1$, and so $U=V$, by Corollary 9.
Since $Aut(V/I(V))$ is abelian, $V/I(V)$ is central in $S/I(V)$,
and so $X/I(V)$ is a central extension of $M_{7,3}$ 
by a torsion free abelian group of rank 1.
Since $H^2(M_{7,3};\mathbb{Z})\cong\mathbb{Z}/3\mathbb{Z}$, 
this extension splits over $C_7$.
Hence $X/I(V)$ has a subgroup of order 7.
This subgroup of order 7 must act trivially on $I(V)$,
since $S$ is torsion free and $I(V)$ is abelian of rank 6.
But then $I(V)$ is central in $S$, and $S$ is virtually nilpotent.
This contradicts our earlier work, and so we cannot have $h(S)<14$.
\end{proof}

\end{document}
